\section*{第 \ref{ch:g12:age-of-the-earth} 章——放射性定年与地球年龄}

\begin{solution}{exo:g12:age-of-the-earth:1}
(a) \omterm{def:g12:age-of-the-earth:closed-system}{封闭系统}；(b) \omterm{def:g12:age-of-the-earth:isochron}{等时线}；(c) \omterm{def:g12:age-of-the-earth:age}{地球年龄}。
\end{solution}

\begin{solution}{exo:g12:age-of-the-earth:2}
斜长石（0.35）、角闪石（1.20）、全岩（3.10）、钾长石（6.40），
黑云母（24.0）。黑云母的数据点最高：它的铷-87 相对于锶-86 最多，
因而其中的铷自
花岗岩形成以来生成的锶-87 最多。
\end{solution}

\begin{solution}{exo:g12:age-of-the-earth:3}
$T = \ln 2/\lambda = 0.6931/\num{1.3972e-11} \approx \num{4.96e10}$ 年，
即 496 亿年，约为\omterm{def:g12:age-of-the-earth:age}{地球年龄}的 $49.6/4.54 \approx 11$ 倍。
\end{solution}

\begin{solution}{exo:g12:age-of-the-earth:4}
$\lambda t = \num{1.3972e-11} \times \num{3.10e8} \approx 0.004331$，且
$e^{0.004331} - 1 \approx 0.004341$。对于 4.8 亿年，$\lambda t
\approx 0.006707$，$e^{\lambda t} - 1 \approx 0.006729$。两者都接近
$\lambda t$ 本身，因为 $\lambda t$ 很小。
\end{solution}

\begin{solution}{exo:g12:age-of-the-earth:5}
菲利普斯（1860）、开尔文（1862）、乔利（1899）、霍姆斯（1927）、帕特森
（1956）。$4540/100 \approx 45$ 倍。
\end{solution}

\begin{solution}{exo:g12:age-of-the-earth:6}
$0.7080 + 10.0 \times 0.004341 \approx 0.7514$ 为今天的值；
花岗岩形成时则为 0.7080，与花岗岩中的每种矿物相同。
\end{solution}

\begin{solution}{exo:g12:age-of-the-earth:7}
（1）$m = (0.77525 - 0.71806)/(9.40 - 0.90) = 0.05719/8.50 \approx 0.006728$；
$t = \ln(1.006728)/\num{1.3972e-11} \approx \num{4.80e8}$ 年，即 4.8
亿年。（2）$0.71806 + 1.70 \times 0.006728 \approx 0.72950$ 与 $0.71806 + 4.20
\times 0.006728 \approx 0.74632$：两点都在该直线上。（3）是的：花岗岩切割
片岩，因此花岗岩较年轻，且 3.1 亿年小于 4.8 亿年。
\end{solution}

\begin{solution}{exo:g12:age-of-the-earth:8}
若只用一种矿物，其各点都会落在同一处：一个点既无法确定直线，
也无法确定斜率。黑云母远在其他各点的右侧，任一点上的小误差，
都会使过该点与斜长石的连线斜率发生改变，
且影响远大于两点相距很近的情形。
\end{solution}

\begin{solution}{exo:g12:age-of-the-earth:9}
（1）其铷的含量未变，故其 $x$ 不变；锶-87 减少，
因此 $y$ 下降：数据点竖直向下移动。（2）过该点与斜长石的直线
斜率变小，所得年龄偏小。
（3）黑云母不再位于其他矿物所连的直线上：各点不再共线，
这说明某个样品未保持\omterm{def:g12:age-of-the-earth:closed-system}{封闭系统}。
\end{solution}

\begin{solution}{exo:g12:age-of-the-earth:10}
（1）$4540/96 \approx 47$ 倍。（2）若有部分盐离开海洋，例如
被埋藏于沉积物中，海洋所含的盐就少于河流带入的全部盐量，
用现存的盐量除以每年的输入量，所得时间会
偏短。
\end{solution}

\begin{solution}{exo:g12:age-of-the-earth:11}
$m = (34.86 - 10.34)/(50.28 - 9.46) = 24.52/40.82 \approx 0.601$，介于
对应于 45.5 亿年（0.590）与 46.0 亿年（0.610）的斜率之间，约为 45.8
亿年。
\end{solution}

\begin{solution}{exo:g12:age-of-the-earth:12}
（1）玄武岩比值较低，来自地幔；花岗岩比值较高，
至少部分来自熔融的地壳。（2）地幔熔融时，铷进入熔体，
因而富集于地壳中，地壳
随时间生成了更多的锶-87。（3）否：同一岩浆结晶出的矿物共享其初始比值，
而这两种岩石的比值并不相同。
\end{solution}

\begin{solution}{exo:g12:age-of-the-earth:13}
（1）铷-87：$2^{-4.54/49.61} \approx 0.94$，约剩 \qty{94}{\percent}，
剩余；铀-238：$2^{-4.54/4.47} \approx 0.49$，约为一半。（2）$m =
e^{\lambda t} - 1 \approx \num{1.3972e-11} \times 10^6 \approx
\num{1.4e-5}$：即使是瓦尔登花岗岩的黑云母（$x = 24$），
只会比初始比值高出 $0.0003$。锶-87 生成得如此之少，
以致斜率难以与零区分：慢时钟适合测定古老的岩石。
\end{solution}

\begin{solution}{exo:g12:age-of-the-earth:14}
（1）$4.438 + 0.6026 \times 19.0 \approx 15.89$；$15.5$ 比其低 $0.39$，
约 \qty{2.4}{\percent}：沉积物样品接近该直线。（2）0.6026
介于 0.590（45.5 亿年）与 0.610（46.0 亿年）之间：约为 45.8 亿年，在
帕特森给出的 44.8 亿至 46.2 亿年范围之内。
\end{solution}

\begin{solution}{exo:g12:age-of-the-earth:15}
（1）此时储存库的净损失仅为 $46.7 - 24.4 = \qty{22.3}{TW}$，
故同一储存库的寿命延长为 $46.7/22.3 \approx 2.1$ 倍。（2）$96 \times 2.1 \approx
200$，即 2 亿年，仍短了二十多倍。（3）\omterm{def:g12:internal-heat:convection}{地幔
对流}，它将热量从深部向上输送，使近表层的梯度
长时间保持较高，远比仅靠传导时持续得久。
\end{solution}

\begin{solution}{pb:g12:age-of-the-earth:1}
\textbf{1.} 黑云母（$x = 24.0$）。

\textbf{2.} 它们由同一岩浆结晶而成，而从同一熔体中结晶的矿物
会携带该熔体的锶及其同位素比值。

\textbf{3.} 锶-86 是稳定的，且不由任何衰变生成：在\omterm{def:g12:age-of-the-earth:closed-system}{封闭系统}中
其数量不会改变，因此它是一把固定的“标尺”。

\textbf{4.} 锶-87 由铷-87 衰变生成；铷-87
减少了，而锶-86 保持不变。

\textbf{5.} 在 $y = 0.7080$ 处的一条水平线：比值相同，$x$ 不同。

\textbf{6.} $(0.81218 - 0.70952)/(24.0 - 0.35) = 0.10266/23.65 \approx
0.004341$.

\textbf{7.} $(0.72146 - 0.70952)/(3.10 - 0.35) = 0.01194/2.75 \approx
0.004342$.

\textbf{8.} 两者一致到 4000 分之一以内。第一个更可靠：
其两点相距甚远，故比值舍入造成的影响较小。

\textbf{9.} $0.70952 - 0.35 \times 0.004341 \approx 0.7080$.

\textbf{10.} $0.7080 + 6.40 \times 0.004341 \approx 0.73578$：
恰好等于实测值。

\textbf{11.} $e^{\lambda t} = 1 + m = 1.004341$.

\textbf{12.} $\lambda t = \ln 1.004341 \approx 0.004331$.

\textbf{13.} $t = 0.004331/\num{1.3972e-11} \approx \num{3.10e8}$ 年。

\textbf{14.} 是的：花岗岩切割了片岩，因此据切割律，
它较年轻，且 3.1 亿年小于
4.8 亿年。

\textbf{15.} 3.1 亿年介于 3.589 亿年与 2.989 亿年之间：
属于石炭纪。

\textbf{16.} 花岗岩结晶时岩浆中的
$^{87}\mathrm{Sr}/^{86}\mathrm{Sr}$ 比值。

\textbf{17.} 花岗岩：$0.7080 > 0.7035$。

\textbf{18.} 地幔熔化时，铷进入熔体，因此铷已
在地壳中富集；铷越多，地壳生成的
锶-87 就越多。

\textbf{19.} 该岩浆至少部分来自熔融的
陆壳，而非仅来自地幔。

\textbf{20.} 瓦尔登花岗岩的年龄为\textbf{3.1 亿年}。
\end{solution}
